% sdl_bgi-quickref.tex
%
% Guido Gonzato, November 15, 2016

\documentclass[a4paper,11pt]{article}
\usepackage{graphicx}
\thispagestyle{empty}
\usepackage{color}
\usepackage{alltt}
\usepackage{framed}
\usepackage[margin=3cm]{geometry}
\usepackage[colorlinks,urlcolor=blue]{hyperref}

\frenchspacing
\setlength{\parindent}{0pt}
\setlength{\parskip}{3pt}

\newcommand{\SDLbgi}{SDL\_bgi}

\newcommand{\version}{2.0.8}        % !!! Change here !!!

\newcommand{\V}{\texttt{void}}      % void
\newcommand{\I}{\texttt{int}}       % int
\newcommand{\U}{\texttt{Uint8}}     % Uint8
\newcommand{\C}{\texttt{char *}}    % char *
\newcommand{\func}[1]{\textbf{#1}}  % function
\newcommand{\A}[1]{\emph{#1}}       % parameter
\newcommand{\F}[1]{\texttt{#1}}     % file
\newcommand{\T}[1]{\texttt{#1}}     % constant

\definecolor{defaultshadecolor}{rgb}{0.85,0.85,1}
\definecolor{shadecolor}{rgb}{0.85,0.85,1}
\definecolor{gray}{rgb}{0.7,0.7,0.7}
\definecolor{charcolor}{rgb}{1,0.85,0.85}
\definecolor{lightred}{rgb}{1,0.7,0.7}
\definecolor{darkblue}{rgb}{0,0,0.5}
\definecolor{darkred}{rgb}{0.75,0,0}

\newenvironment{bgi}
{ % beg def
  \begin{snugshade}
}
{ % end def
  \end{snugshade}
}

% -----

\begin{document}

\title{\SDLbgi{} \version{} Quick Reference}

\author{Guido Gonzato, PhD}

\date{\today}

\maketitle

\begin{center}
  \includegraphics[width=0.5\textwidth]{SDL_bgi_logo.png}  
\end{center}

% -----

\section{Introduction}

\SDLbgi{} is an SDL2-based implementation of Borland Graphics
Interface (BGI), a graphics library that was part of Turbo/Borland C
compilers for DOS. BGI represented the \emph{de facto} standard for
computer graphics in the late eighties--early nineties, especially in
education.

\SDLbgi{} is one of the easiest ways to do graphics programming in C.
It is much simpler (but obviously, less complete) than SDL, OpenGL and
the like. Teachers may find \SDLbgi{} a useful tool for introductory
computer graphics courses.

For example, this is a minimal program that opens a window and draws
1000 random lines:

\begin{small}
\begin{verbatim}
#include <SDL2/SDL_bgi.h>

int main (void)
{
  int i, gd = DETECT, gm;
  initgraph (&gd, &gm, "");
  setbkcolor (BLACK);
  cleardevice ();
  outtextxy (0, 0, "Drawing 1000 lines...");
  for (i = 0; i < 1000; i++) {
    setcolor (1 + random (15));
    line ( random(getmaxx()), random(getmaxy()),
    random(getmaxx()), random(getmaxy()) );
  }
  getch ();
  closegraph ();
  return 0;
}
\end{verbatim}
\end{small}

The program includes the header file \F{SDL\_bgi.h}, which contains
all necessary definitions. The call to \func{initgraph()} opens a
window; from now on, graphics functions may be called.
\func{closegraph()} closes the window.
 
Within the window, pixel coordinates range from (0, 0), the upper left
corner, to (\func{get\-maxx()}, \func{getmaxy()}), the lower right
corner.

Some graphic functions set the coordinates of the last drawing
position, defined as CP (Current Position). At any given moment, a
foreground, background and fill colour, line style, line thickness,
and fill pattern, are defined. A viewport (subwindow) may also be
defined, with or without clipping. All of these parameters can be
changed using appropriate functions.

% -----

\section{Constants}

Many constants are defined in \F{SDL\_bgi.h}. The most important are
the following:

\begin{small}
\begin{verbatim}
#define _SDL_BGI_H
#define SDL_BGI_VERSION 2.0.7

// everything gets drawn here

extern SDL_Window   *bgi_window;
extern SDL_Renderer *bgi_renderer;
extern SDL_Texture  *bgi_texture;

// available visual pages

#define VPAGES 4

// BGI fonts

#define DEFAULT_FONT     0 // 8x8
#define TRIPLEX_FONT     1 // all other fonts are not implemented
#define SMALL_FONT       2
#define SANSSERIF_FONT   3
#define GOTHIC_FONT      4
#define BIG_FONT         5
#define SCRIPT_FONT      6
#define SIMPLEX_FONT     7
#define TRIPLEX_SCR_FONT 8
#define COMPLEX_FONT     9
#define EUROPEAN_FONT    10
#define BOLD_FONT        11

#define HORIZ_DIR      0
#define VERT_DIR       1

#define USER_CHAR_SIZE 0

#define LEFT_TEXT      0
#define CENTER_TEXT    1
#define RIGHT_TEXT     2
#define BOTTOM_TEXT    0
#define TOP_TEXT       2

// BGI colours

#define MAXCOLORS      15
#define BLACK           0
#define BLUE            1
#define GREEN           2
#define CYAN            3
#define RED             4
#define MAGENTA         5
#define BROWN           6
#define LIGHTGRAY       7
#define DARKGRAY        8
#define LIGHTBLUE       9
#define LIGHTGREEN     10
#define LIGHTCYAN      11
#define LIGHTRED       12
#define LIGHTRED       12
#define LIGHTMAGENTA   13
#define YELLOW         14
#define WHITE          15

// line style, thickness, and drawing mode

#define NORM_WIDTH      1
#define THICK_WIDTH     3

#define SOLID_LINE      0
#define DOTTED_LINE     1
#define CENTER_LINE     2
#define DASHED_LINE     3
#define USERBIT_LINE    4

#define COPY_PUT        0
#define XOR_PUT         1
#define OR_PUT          2
#define AND_PUT         3
#define NOT_PUT         4

// fill styles

#define EMPTY_FILL      0
#define SOLID_FILL      1
#define LINE_FILL       2
#define LTSLASH_FILL    3
#define SLASH_FILL      4
#define BKSLASH_FILL    5
#define LTBKSLASH_FILL  6
#define HATCH_FILL      7
#define XHATCH_FILL     8
#define INTERLEAVE_FILL 9
#define WIDE_DOT_FILL   10
#define CLOSE_DOT_FILL  11
#define USER_FILL       12

// mouse buttons

#define WM_LBUTTONDOWN  SDL_BUTTON_LEFT
#define WM_MBUTTONDOWN  SDL_BUTTON_MIDDLE
#define WM_RBUTTONDOWN  SDL_BUTTON_RIGHT
#define WM_WHEELUP      4 // SDL_MOUSEWHEEL no longer exists
#define WM_WHEELDOWN    5 // ditto
#define WM_MOUSEMOVE    SDL_MOUSEMOTION

#define PALETTE_SIZE    4096

#define KEY_HOME        SDLK_HOME
#define KEY_LEFT        SDLK_LEFT
#define KEY_UP          SDLK_UP
#define KEY_RIGHT       SDLK_RIGHT
#define KEY_DOWN        SDLK_DOWN
#define KEY_PGUP        SDLK_PAGEUP
#define KEY_PGDN        SDLK_PAGEDOWN
#define KEY_END         SDLK_END
#define KEY_INSERT      SDLK_INSERT
#define KEY_DELETE      SDLK_DELETE
#define KEY_F1          SDLK_F1
#define KEY_F2          SDLK_F2
#define KEY_F3          SDLK_F3
#define KEY_F4          SDLK_F4
#define KEY_F5          SDLK_F5
#define KEY_F6          SDLK_F6
#define KEY_F7          SDLK_F7
#define KEY_F8          SDLK_F8
#define KEY_F9          SDLK_F9
#define KEY_F10         SDLK_F10
#define KEY_F11         SDLK_F11
#define KEY_F12         SDLK_F12
#define KEY_LEFT_CTRL   SDLK_LCTRL
#define KEY_RIGHT_CTRL  SDLK_RCTRL
#define KEY_LEFT_SHIFT  SDLK_LSHIFT
#define KEY_RIGHT_SHIFT SDLK_RSHIFT
#define KEY_LEFT_ALT    SDLK_LALT
#define KEY_LEFT_WIN    SDLK_LSUPER
#define KEY_RIGHT_WIN   SDLK_RSUPER
#define KEY_ALT_GR      SDLK_MODE
#define KEY_TAB         SDLK_TAB
#define KEY_BS          SDLK_BACKSPACE
#define KEY_RET         SDLK_RETURN
#define KEY_PAUSE       SDLK_PAUSE
#define KEY_SCR_LOCK    SDLK_SCROLLOCK
#define KEY_ESC         SDLK_ESCAPE

// graphics modes

#define DETECT          0
#define grOk            0
#define SDL_CGALO       0
#define SDL_CGAHI       1
#define EGA             2
#define EGAHI           2
#define EGALO           2
#define SDL_EGA         2
#define SDL             3
#define VGA             3
#define SDL_VGA         3
#define SDL_640x480     3
#define SDL_HERC        4
#define SDL_PC3270      5
#define SVGA            6
#define SDL_SVGALO      6
#define SDL_800x600     6
#define SDL_SVGAMED1    7
#define SDL_1024x768    7
#define SDL_SVGAMED2    8
#define SDL_1152x900    8
#define SDL_SVGAHI      9
#define SDL_1280x1024   9
#define SDL_WXGA        10
#define SDL_1366x768    10
#define SDL_USER        11
#define SDL_FULLSCREEN  12
\end{verbatim}
\end{small}

\section{Structs}

Some of the BGI functions use the following structs:

\begin{small}
\begin{verbatim}
struct arccoordstype {
  int x;
  int y;
  int xstart;
  int ystart;
  int xend;
  int yend;
};

struct date {
  int da_year;
  int da_day;
  int da_mon;
};

struct fillsettingstype {
  int pattern;
  int color;
};

struct linesettingstype {
  int linestyle;
  unsigned int upattern;
  int thickness;
};

struct palettetype {
  unsigned char size;
  signed char colors[MAXCOLORS + 1];
};

struct textsettingstype {
  int font;
  int direction;
  int charsize;
  int horiz;
  int vert;
};

struct viewporttype {
  int left;
  int top;
  int right;
  int bottom;
  int clip;
};

struct bgi_info {
  int colour_index;
  char *colour_name;
  unsigned long pixel_value;
};

struct rgb_colour {
  int colour_index;
  unsigned long pixel_value;
};
\end{verbatim}
\end{small}

% -----

\section{Note for SDL2 Programmers}

The following variables, declared in \F{SDL\_bgi.h}, are accessible to
the programmer:

\begin{small}
\begin{verbatim}
SDL_Window   *bgi_window;
SDL_Renderer *bgi_renderer;
SDL_Texture  *bgi_texture;
\end{verbatim}
\end{small}

and can be used by native SDL2 functions. That is, you can use BGI and
native SDL2 functions together, as in the following code snippet:

\begin{small}
\begin{verbatim}
SDL_Surface *bitmap;
SDL_Texture *texture;
...
bitmap = SDL_LoadBMP ("picture.bmp");
texture = SDL_CreateTextureFromSurface (bgi_renderer, bitmap);
SDL_RenderCopy (bgi_renderer, texture, NULL, NULL);
SDL_RenderPresent (bgi_renderer);
...
\end{verbatim}
\end{small}

% -----

\section{Standard BGI Graphics Functions}

The following are standard BGI functions, as implemented for example
in Turbo C. They are all prototyped in \F{SDL\_bgi.h}.

Unless otherwise specified, graphics routines draw shapes using the
current drawing colour, i.e.\ as specified by \func{setcolor()}.

% -----

\begin{bgi}
  \V{} \func{arc} (\I{} \A{x}, \I{} \A{y}, \I{} \A{stangle}, \I{}
  \A{endangle}, \I{} \A{radius});
\end{bgi}

Draws a circular arc centered at (\A{x}, \A{y}), with a radius given
by \A{radius}, traveling from \A{stangle} to \A{endangle}. The angle
for \func{arc()} is measured counterclockwise, with 0 degrees at 3 o'
clock, 90 degrees at 12 o' clock, etc.

\textbf{Note}: The \A{linestyle} parameter does not affect arcs,
circles, ellipses, or pie\-slices. Only the \A{thickness} parameter is
used.

% -----

\begin{bgi}
\V{} \func{bar} (\I{} \A{left}, \I{} \A{top}, \I{} \A{right}, \I{}
\A{bottom});
\end{bgi}

Draws a filled-in rectangle (bar), using the current fill colour and
fill pattern. The bar is not outlined; to draw an outlined
two-dimensional bar, use \func{bar3d()} with \A{depth} equal to 0.

% -----

\begin{bgi}
\V{} \func{bar3d} (\I{} \A{left}, \I{} \A{top}, \I{} \A{right}, \I{}
\A{bottom}, \I{} \A{depth}, \I{} \A{topflag});
\end{bgi}

Draws a three-dimensional, filled-in rectangle (bar), using the
current fill colour and fill pattern. The three-dimensional outline of
the bar is drawn in the current line style and color. The bar's depth,
in pixels, is given by \A{depth}. If \A{topflag} is nonzero, a top is
put on.

% -----

\begin{bgi}
\V{} \func{circle} (\I{} \A{x}, \I{} \A{y}, \I{} \A{radius});
\end{bgi}

Draws a circle of the given \A{radius} at (\A{x}, \A{y}).

\textbf{Note}: The \A{linestyle} parameter does not affect arcs,
circles, ellipses, or pie\-slices. Only the \A{thickness} parameter is
used.

% -----

\begin{bgi}
\V{} \func{cleardevice} (\V{});
\end{bgi}

Clears the graphics screen, filling it with the current background
color. The CP is moved to (0, 0).

% -----

\begin{bgi}
\V{} \func{clearviewport} (\V{});
\end{bgi}

Clears the viewport, filling it with the current background color. The
drawing CP is moved to (0, 0), relative to the viewport.

% -----

\begin{bgi}
\V{} \func{closegraph} (\V{});
\end{bgi}

Closes the graphics system.

% -----

\begin{bgi}
\V{} \func{detectgraph} (\I{} \A{*graphdriver}, \I{} \A{*graphmode});
\end{bgi}

Detects the graphics driver and graphics mode to use. In \SDLbgi, the
default is 800x600 (SVGA).

% -----

\begin{bgi}
\V{} \func{drawpoly} (\I{} \A{numpoints}, \I{} \A{*polypoints});
\end{bgi}

Draws a polygon of \A{numpoints} vertices. \A{polypoints} is a
pointer to a sequence of (2 * \A{numpoints}) integers; each pair gives
the \A{x} and \A{y} coordinate of each vertex.

% -----

\begin{bgi}
\V{} \func{ellipse} (\I{} \A{x}, \I{} \A{y}, \I{} \A{stangle}, \I{}
\A{endangle}, \I{} \A{xradius}, \I{} \A{yradius});
\end{bgi}

Draws an elliptical arc centered at (\A{x}, \A{y}), with axes given by
\A{xradius} and \A{yradius}, traveling from \A{stangle} to
\A{endangle}.

% -----

\begin{bgi}
\V{} \func{fillellipse} (\I{} \A{x}, \I{} \A{y}, \I{} \A{xradius},
\I{} \A{yradius});
\end{bgi}

Draws an ellipse centered at (\A{x}, \A{y}), with axes given by
\A{xradius} and \A{yradius}, and fills it using the current fill
color and fill pattern.

% -----

\begin{bgi}
\V{} \func{fillpoly} (\I{} \A{numpoints}, \I{} \A{*polypoints});
\end{bgi}

Draws a polygon of \A{numpoints} vertices and fills it using the
current fill color.

% -----

\begin{bgi}
\V{} \func{floodfill} (\I{} \A{x}, \I{} \A{y}, \I{} \A{border});
\end{bgi}

Fills an enclosed area, containing the \A{x} and \A{y} points bounded
by the \A{border} color. The area is filled using the current fill
color.

% -----

\begin{bgi}
\I{} \func{getactivepage} (\V{});
\end{bgi}

Returns the active page number.

% -----

\begin{bgi}
\V{} \func{getarccoords} (\texttt{struct arccoordstype} \A{*arccoords});
\end{bgi}

Gets the coordinates of the last call to \func{arc()}, filling the
\A{arccoords} structure.

% -----

\begin{bgi}
\V{} \func{getaspectratio} (\I{} \A{*xasp}, \I{} \A{*yasp});
\end{bgi}

Retrieves the current graphics mode's aspect ratio. In \SDLbgi,
\A{xasp} and \A{yasp} are both 10000 (i.e. the pixels are square).

% -----

\begin{bgi}
\I{} \func{getbkcolor} (\V{});
\end{bgi}

Returns the current background color.

% -----

\begin{bgi}
\I{} \func{getcolor} (\V{});
\end{bgi}

Returns the current drawing (foreground) color.

% -----

\begin{bgi}
\texttt{struct palettetype} \func{*getdefaultpalette} (\V{});
\end{bgi}

Returns the palette definition structure.

% -----

\begin{bgi}
\C{}\func{getdrivername} (\V{});
\end{bgi}

Returns a pointer to a string containing the name of the current
graphics driver.

% -----

\begin{bgi}
\V{} \func{getfillpattern} (\C{} \A{pattern}); 
\end{bgi}

Copies the user-defined fill pattern, as set by setfillpattern, into
the 8-byte area pointed to by \A{pattern}.

% -----

\begin{bgi}
\V{} \func{getfillsettings} (\texttt{struct fillsettingstype}
\A{*fillinfo});
\end{bgi}

Fills the \texttt{fillsettingstype} structure pointed to by
\A{fillinfo} with information about the current fill pattern and fill
color.

% -----

\begin{bgi}
\I{} \func{getgraphmode} (\V{});
\end{bgi}

Returns the current graphics mode.

% -----

\begin{bgi}
\V{} \func{getimage} (\I{} \A{left}, \I{} \A{top}, \I{} \A{right},
\I{} \A{bottom}, \V{} \A{*bitmap});
\end{bgi}

Copies a bit image of the specified region into the memory pointed by
\A{bitmap}.

% -----

\begin{bgi}
\V{} \func{getlinesettings} (\texttt{struct linesettingstype}
\A{*lineinfo});
\end{bgi}

Fills the \texttt{linesettingstype} structure pointed by \A{lineinfo}
with information about the current line style, pattern, and thickness.

% -----

\begin{bgi}
\I{} \func{getmaxcolor} (\V{});
\end{bgi}

Returns the maximum color value available (\T{MAXCOLORS}). If RGB
colors are being used, it returns \T{PALETTE\_SIZE}.

% -----

\begin{bgi}
\I{} \func{getmaxmode} (\V{}); 
\end{bgi}

Returns the maximum mode number for the current driver. In \SDLbgi,
the default is \T{SDL\_FULLSCREEN}.

% -----

\begin{bgi}
\I{} \func{getmaxx} (\V{});
\end{bgi}

Returns the maximum \A{x} screen coordinate.

% -----

\begin{bgi}
\I{} \func{getmaxy} (\V{});
\end{bgi}

Returns the maximum \A{y} screen coordinate.

% -----

\begin{bgi}
\texttt{char*} \func{getmodename} (\I{} \A{mode\_number});
\end{bgi}

Returns a pointer to a string containing the name of the specified
graphics mode.

% -----

\begin{bgi}
\V{} \func{getmoderange} (\I{} \A{graphdriver}, \I{} \A{*lomode}, \I{}
\A{*himode});
\end{bgi}

Gets the range of valid graphics modes. The \A{graphdriver} parameter
is ignored.

% -----

\begin{bgi}
\V{} \func{getpalette} (\texttt{struct palettetype} \A{*palette});
\end{bgi}

Fills the \texttt{palettetype} structure pointed by \A{palette} with
information about the current palette's size and colors.

% -----

\begin{bgi}
\I{} \func{getpalettesize} (\V{}); 
\end{bgi}

Returns the size of the palette (\T{MAXCOLORS} + 1 or
\T{MAXRGB\-COLORS} + 1).

% -----

\begin{bgi}
\I{} \func{getpixel} (\I{} \A{x}, \I{} \A{y});
\end{bgi}

Returns the color of the pixel located at (\A{x}, \A{y}).

% -----

\begin{bgi}
\V{} \func{gettextsettings} (\texttt{struct textsettingstype}
\A{*texttypeinfo});
\end{bgi}

Fills the \texttt{textsettingstype} structure pointed to by
\A{texttypeinfo} with information about the current text font,
direction, size, and justification.

% -----

\begin{bgi}
\V{} \func{getviewsettings} (\texttt{struct viewporttype} \A{*viewport});
\end{bgi}

Fills the \texttt{viewporttype} structure pointed to by \A{viewport}
with information about the current viewport.

% -----

\begin{bgi}
\I{} \func{getvisualpage} (\V{});
\end{bgi}

Returns the visual page number.

% -----

\begin{bgi}
\I{} \func{getx} (\V{});
\end{bgi}

Returns the current viewport's \A{x} coordinate.

% -----

\begin{bgi}
\I{} \func{gety} (\V{});
\end{bgi}

Returns the current viewport's \A{y} coordinate.

% -----

\begin{bgi}
\V{} \func{graphdefaults} (\V{});
\end{bgi}

Resets all graphics settings to their defaults: sets the viewport to
the entire screen, moves the CP to (0, 0), sets the default palette
colors, the default drawing and background color, the default fill
style and pattern, the default text font and justification.

% -----

\begin{bgi}
\texttt{char*} \func{grapherrormsg} (\I{} \A{errorcode});
\end{bgi}

Returns a pointer to the error message string associated with
\A{errorcode}, returned by \func{graphresult()}.

% -----

\begin{bgi}
\I{} \func{graphresult} (\V{});
\end{bgi}

Returns the error code for the last unsuccessful graphics operation
and resets the error level to grOk.

% -----

\begin{bgi}
\texttt{unsigned} \func{imagesize} (\I{} \A{left}, \I{} \A{top}, \I{}
\A{right}, \I{} \A{bottom});
\end{bgi}

Returns the size in bytes of the memory area required to store a bit
image.

% -----

\begin{bgi}
\V{} \func{initgraph} (\I{} \A{*graphdriver}, \I{} \A{*graphmode},
\C{}\A{pathtodriver});
\end{bgi}

Initializes the graphics system. In \SDLbgi, you can use SDL as
\A{graphdriver}, then choose a suitable graphics mode (listed in
\F{graphics.h}) as \A{graphmode}. The \A{pathtodriver} argument is
ignored; you can also use NULL for \A{*graphdriver} and
\A{*graphmode}.

After \func{initgraph()}, all graphics commands are immediately
displayed, as in the original BGI. This makes drawing very slow; you
may want to use \func{initwindow()} instead.

% -----

\begin{bgi}
\V{} \func{initwindow} (\I{} \A{width}, \I{} \A{height});
\end{bgi}

Initializes the graphics system, opening a \A{width}$\times$\A{height}
window. If either \A{width} or \A{height} is 0, then
SDL\-\_FULL\-SCREEN will be used.

The user is in charge of updating the screen using \func{refresh()}.

% -----

\begin{bgi}
\I{} \func{installuserdriver} (\C{}\A{name}, \I{} (\A{*detect})(\V{})); 
\end{bgi}

Unimplemented; not used by \SDLbgi.

% -----

\begin{bgi}
\I{} \func{installuserfont} (\C{}\A{name}); 
\end{bgi}

Unimplemented; not used by \SDLbgi.

% -----

\begin{bgi}
\V{} \func{line} (\I{} \A{x1}, \I{} \A{y1}, \I{} \A{x2}, \I{} \A{y2});
\end{bgi}

Draws a line between two specified points; the CP is not updated.

% -----

\begin{bgi}
\V{} \func{linerel} (\I{} \A{dx}, \I{} \A{dy});
\end{bgi}

Draws a line from the CP to a point that is (\A{dx},\A{dy}) pixels
from the CP. The CP is then advanced by (\A{dx}, \A{dy}).

% -----

\begin{bgi}
\V{} \func{lineto} (\I{} \A{x}, \I{} \A{y});
\end{bgi}

Draws a line from the CP to (\A{x}, \A{y}), then moves the CP to
(\A{dx}, \A{dy}).

% -----

\begin{bgi}
\V{} \func{moverel} (\I{} \A{dx}, \I{} \A{dy});
\end{bgi}

Moves the CP by (\A{dx}, \A{dy}) pixels.

% -----

\begin{bgi}
\V{} \func{moveto} (\I{} \A{x}, \I{} \A{y});
\end{bgi}

Moves the CP to the position (\A{x}, \A{y}), relative to the
viewport.

% -----

\begin{bgi}
\V{} \func{outtext} (\C{}\A{textstring});
\end{bgi}

Outputs \A{textstring} at the CP.

% -----

\begin{bgi}
\V{} \func{outtextxy} (\I{} \A{x}, \I{} \A{y}, \C{}\A{textstring});
\end{bgi}

Outputs \A{textstring} at (\A{x}, \A{y}).

% -----

\begin{bgi}
\V{} \func{pieslice} (\I{} \A{x}, \I{} \A{y}, \I{} \A{stangle}, \I{}
\A{endangle}, \I{} \A{radius});
\end{bgi}

Draws and fills a pie slice centered at (\A{x}, \A{y}), with a radius
given by \A{radius}, traveling from \A{stangle} to \A{endangle}. The
pie slice is filled using the current fill colour.

% -----

\begin{bgi}
\V{} \func{putimage} (\I{} \A{left}, \I{} \A{top}, \V{} \A{*bitmap},
\I{} \A{op});
\end{bgi}

Puts the bit image pointed to by \A{bitmap} onto the screen, with the
upper left corner of the image placed at (\A{left}, \A{top}). \A{op}
specifies the drawing mode (COPY\_PUT, etc).

% -----

\begin{bgi}
\V{} \func{putpixel} (\I{} \A{x}, \I{} \A{y}, \I{} \A{color}); 
\end{bgi}

Plots a point at (\A{x},\A{y}) in the color defined by \A{color}.

% -----

\begin{bgi}
\V{} \func{rectangle} (\I{} \A{left}, \I{} \A{top}, \I{} \A{right},
\I{} \A{bottom});
\end{bgi}

Draws a rectangle delimited by (\A{left},\A{top}) and
(\A{right},\A{bottom}).

% -----

\begin{bgi}
\V{} \func{refresh} (\V{});
\end{bgi}

Updates the screen.

% -----

\begin{bgi}
\I{} \func{registerbgidriver} (\V{} (\A{*driver})(\V{})); 
\end{bgi}

Unimplemented; not used by \SDLbgi.

% -----

\begin{bgi}
\I{} \func{registerbgifont} (\V{} (\A{*font})(\V{})); 
\end{bgi}

Unimplemented; not used by \SDLbgi.

% -----

\begin{bgi}
\V{} \func{restorecrtmode} (\V{}); 
\end{bgi}

Hides the graphics window.

% -----

\begin{bgi}
\V{} \func{sector} (\I{} \A{x}, \I{} \A{y}, \I{} \A{stangle}, \I{}
\A{endangle}, \I{} \A{xradius}, \I{} \A{yradius});
\end{bgi}

Draws and fills an elliptical pie slice centered at (\A{x}, \A{y}),
horizontal and vertical radii given by \A{xradius} and \A{yradius},
traveling from \A{stangle} to \A{endangle}.

% -----

\begin{bgi}
\V{} \func{setactivepage} (\I{} \A{page});
\end{bgi}

Makes \A{page} the active page for all subsequent graphics output.

% -----

\begin{bgi}
\V{} \func{setallpalette} (\texttt{struct palettetype *}\A{palette}); 
\end{bgi}

Sets the current palette to the values given in \A{palette}.

% -----

\begin{bgi}
\V{} \func{setaspectratio} (\I{} \A{xasp}, \I{} \A{yasp});
\end{bgi}

Changes the default aspect ratio of the graphics. In \SDLbgi, this
function is not necessary since the pixels are square.

% -----

\begin{bgi}
\V{} \func{setbkcolor} (\I{} \A{color});
\end{bgi}

Sets the current background color using the default palette.

% -----

\begin{bgi}
\V{} \func{setcolor} (\I{} \A{color});
\end{bgi}

Sets the current drawing color using the default palette.

% -----

\begin{bgi}
\V{} \func{setfillpattern} (\C{}\A{upattern}, \I{} \A{color}); 
\end{bgi}

Sets a user-defined fill pattern. \A{upattern} is a pointer to a
sequence of 8 bytes; each byte corresponds to 8 pixels in the pattern;
each bit set to 1 is plotted as a pixel.

% -----

\begin{bgi}
\V{} \func{setfillstyle} (\I{} \A{upattern}, \I{} \A{color}); 
\end{bgi}

Sets the fill pattern and fill color. \A{upattern} is a pointer to a
sequence of 8 bytes, with each byte corresponding to 8 pixels in the
pattern.

% -----

\begin{bgi}
unsigned \func{setgraphbufsize} (\texttt{unsigned} \A{bufsize}); 
\end{bgi}

Unimplemented; not used by \SDLbgi.

% -----

\begin{bgi}
\V{} \func{setgraphmode} (\I{} \A{mode}); 
\end{bgi}

Shows the window that was hidden by \func{restorecrtmode()}. The
\A{mode} parameter is ignored

% -----

\begin{bgi}
\V{} \func{setlinestyle} (\I{} \A{linestyle}, \texttt{unsigned}
\A{upattern}, \I{} \A{thickness});
\end{bgi}

Sets the line width and style for all lines drawn by \func{line()},
\func{lineto()}, \func{rectangle()}, \func{drawpoly()}, etc. The line
style can be \T{SOLID\_LINE}, \T{DOTTED\_LINE}, \T{CENTER\_LINE},
\T{DASHED\_LINE}, or \T{USERBIT\_LINE}; in the latter case, the user
provides a 16-bit number (\A{upattern}) whose bits set to 1 will be
plotted as pixels.

The line thickness can be set with \T{NORM\_WIDTH} or
\T{THICK\-\_WIDTH}.

Arcs, circles, ellipses, and pieslices are not affected by
\A{linestyle}, but are affected by \A{thickness}.

% -----

\begin{bgi}
\V{} \func{setpalette} (\I{} \A{colornum}, \I{} \A{color}); 
\end{bgi}

Changes the standard palette \A{colornum} to \A{color}.

% -----

\begin{bgi}
\V{} \func{settextjustify} (\I{} \A{horiz}, \I{} \A{vert});
\end{bgi}

Sets text justification. Text output will be justified around the CP
horizontally and vertically; settings are \T{LEFT\_TEXT},
\T{CENTER\_TEXT}, \T{RIGHT\-\_TEXT}, \T{BOTTOM\_TEXT}, and
\T{TOP\_TEXT}.

% -----

\begin{bgi}
\V{} \func{settextstyle} (\I{} \A{font}, \I{} \A{direction}, \I{}
\A{charsize});
\end{bgi}

Sets the text font (only \T{DEFAULT\_FONT} is actually available), the
direction in which text is displayed (\T{HORIZ\_DIR}, \T{VERT\_DIR}),
and the size of the characters. If \A{charsize} is an integer number,
the text will be scaled by that number; if it is 0, the text will be
scaled by \func{setusercharsize()}.

% FIX IT!

% -----

\begin{bgi}
\V{} \func{setusercharsize} (\I{} \A{multx}, \I{} \A{divx}, \I{}
\A{multy}, \I{} \A{divy});
\end{bgi}

Lets the user change the character width and height. If a previous
call to \func{settextstyle()} set \A{charsize} to 0, the default width
is scaled by \A{multx}/\A{divx}, and the default height is scaled by
\A{multy}/\A{divy}.

% -----

\begin{bgi}
\V{} \func{setviewport} (\I{} \A{left}, \I{} \A{top}, \I{} \A{right},
\I{} \A{bottom}, \I{} \A{clip});
\end{bgi}

Sets the current viewport for graphics output. If \A{clip} is nonzero,
all drawings will be clipped (truncated) to the current viewport.

% -----

\begin{bgi}
\V{} \func{setvisualpage} (\I{} \A{page});
\end{bgi}

Sets the visual graphics page number.

% -----

\begin{bgi}
\V{} \func{setwritemode} (\I{} \A{mode});
\end{bgi}

Sets the writing mode for line drawing. \A{mode} can be \T{COPY\_PUT},
\T{XOR\-\_PUT}, \T{OR\_PUT}, \T{AND\_PUT}, and \T{NOT\_PUT}.

% -----

\begin{bgi}
\I{} \func{textheight} (\C{}\A{textstring});
\end{bgi}

Returns the height in pixels of a string.

% -----

\begin{bgi}
\I{} \func{textwidth} (\C{}\A{textstring});
\end{bgi}

Returns the height in pixels of a string.

% -----

\section{Non-Graphics Functions and Macros}

\begin{bgi}
\V{} \func{delay} (\I{} \A{millisec});
\end{bgi}

Waits for \A{millisec} milliseconds. In ``slow mode'', a screen
refresh is performed.

% -----

\begin{bgi}
\I{} \func{getch} (\V{});
\end{bgi}

Waits for a key and returns its ASCII code. In ``slow mode'', a screen
refresh is performed.

% -----

\begin{bgi}
\I{} \func{kbhit} (\V{});
\end{bgi}

Returns 1 when a key is pressed. In ``slow mode'', a screen refresh is
performed.

% -----

\begin{bgi}
\I{} \func{random} (\I{} \A{range}) (macro)
\end{bgi}

Returns a random number between 0 and \A{range - 1}.

% ----- ----- -----
	
\section{\SDLbgi{} Additions}

\begin{bgi}
\I{} \func{COLOR} (\I{} \A{r}, \I{} \A{g}, \I{} \A{b});
\end{bgi}

Can be used as an argument for \func{setcolor()} and
\func{setbkcolor()} to set an RBG color.

% -----

\begin{bgi}
\I{} RED\_VALUE (\I{} \A{color})
\end{bgi}

Returns the red component of an RGB color.

% -----

\begin{bgi}
\I{} GREEN\_VALUE (\I{} \A{color})
\end{bgi}

Returns the green component of an RGB color.

% -----

\begin{bgi}
\I{} BLUE\_VALUE (\I{} \A{color})
\end{bgi}

Returns the blue component of an RGB color.

% -----

\begin{bgi}
\I{} IS\_BGI\_COLOR (\I{} \A{color});
\end{bgi}

Returns 1 if the current color is a standard BGI color (not RGB). The
argument is actually redundant.

% -----

\begin{bgi}
\I{} IS\_RGB\_COLOR (\I{} \A{color});
\end{bgi}

Returns 1 if the current color is in RGB mode. The argument is
actually redundant.

% -----

\begin{bgi}
\V{} \func{event} (\V{});
\end{bgi}

Returns 1 if an event (mouse click or key press) has occurred.

% -----

\begin{bgi}
\V{} \func{eventtype} (\V{});
\end{bgi}

Returns the type of the last event; either \T{SDL\_KEYPRESS} or
\T{SDL\_MOUSEBUTTONDOWN}.

% -----

\begin{bgi}
\I{} \func{getevent} (\V{});
\end{bgi}

Waits for a keypress or mouse click, and returns the code of the mouse
button or key that was pressed.

% -----

\begin{bgi}
\I{} \func{getmaxheight} (\V{});
\end{bgi}

Returns the screen (root window) height.

% -----

\begin{bgi}
\I{} \func{getmaxwidth} (\V{});
\end{bgi}

Returns the screen (root window) width.

% -----

\begin{bgi}
\V{} \func{getmouseclick} (\I{} \A{kind}, \I{}\A{*x}, \I{}\A{*y});
\end{bgi}

Sets the \A{x},\A{y} coordinates of the last \A{kind} button click
expected by \func{ismouseclick()}.

% -----

\begin{bgi}
\I{} \func{initwindow} (\I{} \A{width}, \I{} \A{height});
\end{bgi}

Opens a \A{width}$\times$\A{height} graphics window.

% -----

\begin{bgi}
\I{} \func{ismouseclick} (\I{} \A{kind});
\end{bgi}

Returns 1 if the \A{kind} mouse button was clicked.

% -----

\begin{bgi}
\I{} \func{mouseclick} (\V{});
\end{bgi}

Returns the code of the mouse button that was clicked, or 0 if none
was clicked.

% -----

\begin{bgi}
\I{} \func{mousex} (\V{});
\end{bgi}

Returns the X coordinate of the last mouse click.

% -----

\begin{bgi}
\I{} \func{mousey} (\V{});
\end{bgi}

Returns the Y coordinate of the last mouse click.

% -----

\begin{bgi}
\V{} \func{\_putpixel} (\I{} \A{x}, \I{} \A{y}); 
\end{bgi}

Plots a point at (\A{x},\A{y}) using the current drawing color. This
function is faster than \func{putpixel()}, since it is not immediately
displayed.

% -----

\begin{bgi}
\V{} \func{readimagefile} (\C{} \A{filename}, \I{} \A{x1}, \I{} \A{y1}, 
\I{} \A{x2}, \I{} \A{y2});
\end{bgi}

Reads a \F{.bmp} file and displays it immediately at (\A{x1}, \A{y1}).
If (\A{x2}, \A{y2}) are not 0, the bitmap is stretched to fit the
rectangle \A{x1},\A{y1}---\A{x2},\A{y2}; otherwise, the bitmap is
clipped as necessary.

% -----

\begin{bgi}
\V{} \func{refresh} (\V{}); 
\end{bgi}

Updates the screen contents, i.e.\ displays all graphics.

% -----

\begin{bgi}
\V{} \func{sdlbgifast} (\V{}); 
\end{bgi}

Triggers ``fast mode'', i.e.\ \func{refresh()} is needed to display
graphics.

% -----

\begin{bgi}
\V{} \func{sdlbgislow} (\V{}); 
\end{bgi}

Triggers ``slow mode'', i.e.\ \func{refresh()} is not needed to
display graphics.

% -----

\begin{bgi}
\V{} \func{setalpha} (\I{} \A{col}, \U{} \A{alpha}); 
\end{bgi}

Sets alpha transparency for colour \A{col} to \A{alpha} (0--255).

% -----


\begin{bgi}
\V{} \func{setbkrgbcolor} (\I{} \A{n}); 
\end{bgi}

Sets the current background color using using the \A{n}-th color index
in the RGB palette.

% -----

\begin{bgi}
\V{} \func{setrgbcolor} (\I{} \A{n}); 
\end{bgi}

Sets the current drawing color using the \A{n}-th color index in the
RGB palette.

% -----

\begin{bgi}
\V{} \func{setrgbpalette} (\I{} \A{n}, \I{} \A{r}, \I{} \A{g}, \I{} \A{b}); 
\end{bgi}

Sets the \A{n}-th entry in the RGB palette specifying the \A{r},
\A{g}, and \A{b} components.

Using \func{setrgbpalette()} and \func{setrgbcolor()} is faster than
setting colors with \func{setcolor()} with a \func{COLOR()} argument.

% -----

\begin{bgi}
\I{} \func{swapbuffers} (\V{});
\end{bgi}

Swaps the current active and the current visual graphics pages.

% -----

\begin{bgi}
\I{} \func{xkbhit} (\V{});
\end{bgi}

Returns 1 when any key is pressed, including Ctrl, Shift, etc. In
``slow mode'', a screen refresh is performed.

% -----

\begin{bgi}
\V{} \func{writeimagefile} (\C{} \A{filename}, \I{} \A{left}, \I{}
\A{top}, \I{} \A{right}, \I{} \A{bottom});
\end{bgi}

Writes a \F{.bmp} file from the screen rectangle defined by
\A{left},\A{top}---\A{right},\A{bottom}.

% -----

\bigskip

\rule{\linewidth}{0.2pt}

This document is a free manual, released under the GNU Free
Documentation License (FDL) v. 1.3 or later.

\end{document}

% ----- end of file sdl_bgi-quickref.tex
